\section{Block Parallelism}
\label{sec:method}

\subsection{Limitations of conventional context parallelism}
\label{sec:cp-limitations}

Conventional CP is poorly matched to BDLM training because it partitions the
entire length-$2L$ clean-plus-corrupted representation by token position and
performs distributed attention across those position shards. Both clean and
corrupted K/V therefore participate in cross-rank attention communication.
The corrupted K/V for block $b$ are used only by queries from that block. CP
exchanges these block-specific K/V during forward and returns their gradients
to the position-shard owners during backward, adding communication at every
layer for data used by only one block.

Clean K/V are shared across later block losses. When
$\mathcal{L}_b$ uses a clean K/V shard, attention backward produces a gradient
contribution for that shard. This contribution is accumulated on the owning
rank, where backward continues through the clean projections and earlier
layers. Conventional CP therefore combines the required clean-prefix
communication with block-specific corrupted K/V and gradient communication at
every layer.

\subsection{Block-parallel execution}
\label{sec:block-parallel-training}

Our first insight is that we can eliminate this block-specific communication by
exploiting the BDLM objective's per-block loss separability to expose
target-block computation as a new distributed parallelism dimension. The
corrupted activations of block $b$ contribute only to $\mathcal{L}_b$, while
the clean representations of its observed prefix can contribute to multiple
block losses. BP assigns each complete corrupted-block computation to one rank.
The assigned rank executes the corrupted forward pass, loss, and backward pass
for block $b$ while keeping the block's corrupted Q/K/V, activations, and
gradients local. Different ranks execute different block computations
concurrently (Figure~\ref{fig:block-parallel-overview}(b)).

Let $\mathcal{B}_r$ denote the blocks assigned to rank $r$. Rank $r$ computes
\begin{equation}
    \mathcal{L}_r
    =
    \sum_{b\in\mathcal{B}_r}
    \mathcal{L}_b.
    \label{eq:local-block-loss}
\end{equation}
Linearity of differentiation gives the global gradient identity
\begin{equation}
    \sum_{r=1}^{P}\nabla_\theta\mathcal{L}_r
    =
    \sum_{b=1}^{B}\nabla_\theta\mathcal{L}_b
    =
    \nabla_\theta\mathcal{L}_{\mathrm{BDLM}}.
    \label{eq:block-gradient-sum}
\end{equation}

\subsection{Context-sharded block parallelism}
\label{sec:context-sharded-block-parallelism}

BP keeps corrupted-block computations local, but BP alone replicates clean
prefixes across ranks. A rank can reuse the longest prefix required by its
assigned blocks, but these prefixes overlap across ranks, and ranks assigned
later blocks may compute and store nearly the full clean sequence. Clean
computation and activations can therefore be replicated across up to
$\min(B,P)$ active ranks.

Our second insight is that we can scale BP without this replication because BP
and CP operate on complementary parts of the computation over the same $P$
ranks. BP assigns complete corrupted-block computations while CP shards the
clean-sequence computation shared by those blocks. We call this
context-sharded block parallelism (CSBP). Each rank executes a subset of
complete corrupted-block computations and stores one clean-sequence shard
(Figure~\ref{fig:csbp}). Context parallelism supplies the distributed clean
K/V and accumulates their backward contributions, while BP keeps each
corrupted-block computation on its owner.

\begin{figure}[t]
    \centering
    \resizebox{\linewidth}{!}{\begin{tikzpicture}[
    x=1cm,
    y=1cm,
    font=\sffamily\scriptsize,
    >={Latex[length=2.2mm,width=1.5mm]},
    device/.style={draw=black!58, fill=black!2, rounded corners=2.3pt,
                   line width=0.7pt, minimum width=2.64cm,
                   minimum height=2.16cm},
    ownerdevice/.style={draw=BPOrange, fill=BPOrange!4,
                       rounded corners=2.3pt, line width=1.05pt,
                       minimum width=2.64cm, minimum height=2.16cm},
    clean/.style={draw=BPBlue, fill=BPBlue!13, rounded corners=1.4pt,
                  line width=0.7pt, align=center},
    corrupt/.style={draw=BPOrange, fill=BPOrange!16, rounded corners=1.4pt,
                    dash pattern=on 2.4pt off 1.2pt, line width=0.75pt,
                    align=center},
    shard/.style={minimum width=1.10cm, minimum height=0.58cm,
                  text width=1.00cm, inner xsep=0pt, align=center},
    blockjob/.style={minimum width=2.16cm, minimum height=0.78cm,
                     text width=2.06cm, inner xsep=0pt, align=center},
    callout/.style={draw=black!38, fill=black!3, rounded corners=2.3pt,
                    line width=0.65pt, text width=2.90cm,
                    minimum width=3.18cm, minimum height=1.14cm,
                    inner xsep=0.14cm, inner ysep=0.10cm,
                    align=center},
    fwdclean/.style={draw=BPBlue, line width=1.0pt,
                     -{Latex[length=2.4mm,width=1.6mm]}},
    fwdcorrupt/.style={draw=BPOrange, line width=1.0pt,
                       -{Latex[length=2.4mm,width=1.6mm]}},
    bwdclean/.style={draw=BPBlue, line width=0.95pt,
                     dash pattern=on 3pt off 1.6pt,
                     -{Latex[length=2.4mm,width=1.6mm]}},
    bwdcorrupt/.style={draw=BPOrange, line width=0.95pt,
                       dash pattern=on 3pt off 1.6pt,
                       -{Latex[length=2.4mm,width=1.6mm]}},
    groupclean/.style={draw=BPBlue, line width=1.0pt,
                       {Latex[length=2.4mm,width=1.6mm]}-{Latex[length=2.4mm,width=1.6mm]}},
    groupcorrupt/.style={draw=BPOrange, line width=1.0pt,
                         {Latex[length=2.4mm,width=1.6mm]}-{Latex[length=2.4mm,width=1.6mm]}},
    groupgradclean/.style={draw=BPBlue, line width=0.95pt,
                           dash pattern=on 3pt off 1.6pt,
                           {Latex[length=2.4mm,width=1.6mm]}-{Latex[length=2.4mm,width=1.6mm]}},
    groupgradcorrupt/.style={draw=BPOrange, line width=0.95pt,
                             dash pattern=on 3pt off 1.6pt,
                             {Latex[length=2.4mm,width=1.6mm]}-{Latex[length=2.4mm,width=1.6mm]}},
    connector/.style={draw=black!30, line width=0.55pt},
    cleanconnector/.style={draw=BPBlue!65, line width=0.6pt}
]
\definecolor{BPBlue}{HTML}{356F9F}
\definecolor{BPOrange}{HTML}{B66718}

% Legend.
\node[clean, minimum width=0.34cm, minimum height=0.20cm,
      inner sep=0pt] at (4.02,6.02) {};
\node[anchor=west, text=black!72, font=\sffamily\tiny]
      at (4.27,6.02) {Clean tokens};
\node[draw=BPOrange, fill=BPOrange!15, minimum width=0.34cm, minimum height=0.20cm,
      inner sep=0pt] at (6.57,6.02) {};
\node[anchor=west, text=black!72, font=\sffamily\tiny]
      at (6.82,6.02) {Corrupted tokens};

% ---------------------------------------------------------------------------
% Context parallelism: the combined clean-plus-corrupted sequence is sharded.
% ---------------------------------------------------------------------------
\node[anchor=east, font=\sffamily\scriptsize\bfseries, text=black!82]
      at (1.85,4.28) {(a) CP};

\foreach \x/\r/\n in {3.18/{1}/1,6.18/{2}/2,9.18/{3}/3} {
    \node[device] (cpd\n) at (\x,4.28) {};
    \node[font=\sffamily\scriptsize\bfseries, text=black!76]
          at ([yshift=-0.33cm]cpd\n.north) {Rank \r};
    \node[font=\sffamily\tiny, align=center, text=BPBlue!88!black]
          at ({\x-0.57},4.62) {Clean\\$Q/K/V$};
    \node[font=\sffamily\tiny, align=center, text=BPOrange!90!black]
          at ({\x+0.57},4.62) {Corrupted\\$Q/K/V$};
    \foreach \i in {1,2,3} {
        \pgfmathtruncatemacro{\tokenpos}{3*(\n-1)+\i}
        \node[draw=BPBlue, fill=BPBlue!13, line width=0.65pt,
              minimum width=0.30cm, minimum height=0.40cm,
              inner sep=0pt, font=\sffamily\tiny]
              at ({\x-0.57+0.34*(\i-2)},4.16) {\tokenpos};
        \node[draw=BPOrange, fill=BPOrange!15, line width=0.65pt,
              minimum width=0.30cm, minimum height=0.40cm,
              inner sep=0pt, font=\sffamily\tiny]
              at ({\x+0.57+0.34*(\i-2)},4.16) {\tokenpos};
    }
    \node[font=\sffamily\tiny\bfseries, text=black!72, align=center]
          at ([yshift=0.48cm]cpd\n.south) {$Q$, $dQ$, and $O$\\[-1pt]stay local};
}

% The combined sequence is the CP input; neutral group markers avoid implying
% a particular transport.
\node[anchor=east, font=\sffamily\tiny\bfseries, text=black!68]
      at (1.93,5.61) {Forward: $K/V$};
\draw[groupclean] (2.13,5.74) -- (10.23,5.74);
\draw[groupcorrupt] (2.13,5.48) -- (10.23,5.48);
\foreach \x in {3.18,6.18,9.18} {
    \draw[connector] (\x,5.36) -- (\x,5.74);
}

% Gradients for both parts of the combined input return to shard owners.
\node[anchor=east, font=\sffamily\tiny\bfseries, text=black!68]
      at (1.93,2.89) {Backward: $dK/dV$};
\draw[groupgradclean] (2.13,3.02) -- (10.23,3.02);
\draw[groupgradcorrupt] (2.13,2.76) -- (10.23,2.76);
\foreach \x in {3.18,6.18,9.18} {
    \draw[connector] (\x,3.20) -- (\x,2.76);
}

\node[callout, font=\sffamily\tiny] at (12.34,4.28)
      {$Q$, $dQ$, $O$ stay local\\[-1pt]
       \textbf{Cross-rank: Clean + Corrupted}\\[-1pt]
       \textbf{$K/V$ and $dK/dV$}};

% ---------------------------------------------------------------------------
% CSBP: block b is owned by one rank; other blocks run concurrently.
% ---------------------------------------------------------------------------
\node[anchor=east, font=\sffamily\scriptsize\bfseries, text=black!82]
      at (1.85,1.28) {(b) CSBP};

\node[device, minimum height=2.26cm] (fd1) at (3.18,1.23) {};
\node[device, minimum height=2.26cm] (fd2) at (6.18,1.23) {};
\node[ownerdevice, minimum height=2.26cm] (fd3) at (9.18,1.23) {};
\foreach \r/\n in {1/1,2/2,3/3} {
    \node[font=\sffamily\scriptsize\bfseries, text=black!76]
          at ([yshift=-0.33cm]fd\n.north) {Rank \r};
    \node[font=\sffamily\tiny, text=BPBlue!88!black]
          at ([yshift=0.53cm]fd\n.center) {Clean $Q/K/V$};
    \foreach \i in {1,2,3} {
        \pgfmathtruncatemacro{\tokenpos}{3*(\n-1)+\i}
        \pgfmathsetmacro{\tokenshift}{0.34*(\i-2)}
        \node[draw=BPBlue, fill=BPBlue!13, line width=0.65pt,
              minimum width=0.30cm, minimum height=0.40cm,
              inner sep=0pt, font=\sffamily\tiny]
              at ([xshift=\tokenshift cm,yshift=0.20cm]fd\n.center)
              {\tokenpos};
    }
}
\node[corrupt, draw=BPOrange, fill=BPOrange!15, text=black!85,
      blockjob, font=\sffamily\tiny] at ([yshift=-0.59cm]fd1.center)
      {Block $1$: Corrupted\\$Q/K/V + dQ/dK/dV$ local};
\node[corrupt, draw=BPOrange, fill=BPOrange!15, text=black!85,
      blockjob, font=\sffamily\tiny] at ([yshift=-0.59cm]fd2.center)
      {Block $2$: Corrupted\\$Q/K/V + dQ/dK/dV$ local};
\node[corrupt, blockjob, font=\sffamily\tiny\bfseries]
      at ([yshift=-0.59cm]fd3.center)
      {Block $3$: Corrupted\\$Q/K/V + dQ/dK/dV$ local};

% Clean K/V travels to the owner of block 3; corrupted K/V remains local.
\node[anchor=east, font=\sffamily\tiny\bfseries, text=black!68]
      at (1.93,2.56) {Forward: $K/V$};
\draw[fwdclean] (3.18,2.64) -- (9.18,2.64);
\draw[fwdclean] (6.18,2.48) -- (9.18,2.48);
\draw[cleanconnector] (fd1.north) -- (3.18,2.64);
\draw[cleanconnector] (fd2.north) -- (6.18,2.64);
\draw[cleanconnector] (fd3.north) -- (9.18,2.64);

% Clean-KV gradients return to shard owners; block-3 backward remains on Rank 3.
\node[anchor=east, font=\sffamily\tiny\bfseries, text=black!68]
      at (1.93,-0.08) {Backward: $dK/dV$};
\draw[bwdclean] (9.18,-0.13) -- (3.18,-0.13);
\draw[bwdclean] (9.18,0.02) -- (6.18,0.02);
\draw[cleanconnector] (fd1.south) -- (3.18,-0.13);
\draw[cleanconnector] (fd2.south) -- (6.18,0.02);
\draw[cleanconnector] (fd3.south) -- (9.18,-0.13);

\node[callout, draw=BPOrange!78, fill=BPOrange!7,
      font=\sffamily\tiny] at (12.34,1.28)
      {$Q$, $dQ$, $O$ stay local\\[-1pt]
       \textbf{Cross-rank: Clean $K/V$}\\[-1pt]
       \textbf{and Clean $dK/dV$ only}};
\end{tikzpicture}
}
    \caption{Conventional CP shards the combined clean-plus-corrupted sequence
    by token position. Clean and corrupted $K/V$ cross ranks during forward,
    and clean and corrupted $dK/dV$ cross ranks during backward; $Q$, $dQ$, and
    attention outputs remain with their query owners. CSBP shards only the
    shared clean sequence, so only clean $K/V$ and $dK/dV$ cross ranks. Each
    block owner retains its corrupted $Q/K/V$ and $dQ/dK/dV$ locally.}
    \label{fig:csbp}
\end{figure}

\paragraph{Forward pass.}
The rank assigned target block $b$ computes
\begin{equation}
    O_b
    =
    \operatorname{Attn}\!\left(
        Q_b,
        [K_b;K_{<b}^{\mathrm{clean}}],
        [V_b;V_{<b}^{\mathrm{clean}}]
    \right),
    \label{eq:csbp-block-attention}
\end{equation}
where $\operatorname{Attn}$ denotes scaled dot-product attention. The rank
holds $Q_b$, $K_b$, and $V_b$ locally, while CP supplies the sharded
clean-prefix K/V. Online softmax combines the local within-block contribution
with the distributed clean-prefix
contributions~\citep{milakov2018online,dao2022flashattention}.

The target-block computations reuse overlapping prefixes of the clean
sequence. During one context-parallel attention pass, rank $r$ applies each
clean K/V shard supplied by CP to its clean queries and to every
$b\in\mathcal{B}_r$ whose prefix contains tokens from that shard. The
length-$L$ distributed clean K/V stream therefore serves every local block
whose clean prefix includes that shard. Algorithm~\ref{alg:csbp-attention}
gives the complete per-layer computation in the same notation.

\paragraph{Backward pass.}
The rank assigned block $b$ backpropagates $\mathcal{L}_b$, computing the
gradients of $Q_b$, $K_b$, and $V_b$ locally together with contributions to
the clean-prefix K/V gradients. Exact CP backward accumulates the clean K/V
gradient contributions from all clean and assigned corrupted-block queries on
each clean-shard owner, which continues backward through the clean projections
and earlier layers. Corrupted Q/K/V gradients therefore remain on block owners,
while gradients through clean activations are accumulated on clean-shard
owners. Together, these contributions give the global parameter gradient in
Equation~\ref{eq:block-gradient-sum}. Algorithm~\ref{alg:csbp-backward}
gives the corresponding backward pass. Implementations~\ref{impl:csbp-forward}
and~\ref{impl:csbp-backward} give the production forward and backward fast
paths.

Later target blocks and clean-sequence chunks have longer prefixes and require
more attention work. We devise a dual-end strategy that assigns target blocks
$b$ and $B+1-b$ to the same rank. We use the same early--late pairing for
target blocks and zigzag clean-sequence sharding over the same
ranks~\citep{dubey2024llama3}. Appendix~\ref{app:load-balancing} shows that this
common pairing exactly balances target-block and clean causal attention for
equal partitions and handles non-divisible block counts.

\subsection{Complexity and sources of improvement}
\label{sec:complexity-analysis}

CSBP preserves the exact BDLM computation and gradients
(Proposition~\ref{prop:csbp-exactness}) while eliminating cross-rank
communication for corrupted K/V and avoiding replicated clean prefixes.

\paragraph{Computation.}
Conventional CP and CSBP evaluate the same valid attention pairs, so
both require $\Theta(L^2)$ total attention work and $\Theta(L^2/P)$ work per
rank under balanced assignment. CSBP changes where this work is
executed. Each corrupted block is evaluated on one rank, while CP supplies the
shared clean K/V. The gain comes from keeping block-specific computation local
and eliminating its cross-rank communication while preserving the same
asymptotic attention work.

\paragraph{Communication.}
Conventional CP applies distributed attention to the combined length-$2L$
clean-plus-corrupted representation. CSBP distributes only the
length-$L$ clean K/V, while corrupted K/V remain on their block owners. Each
communicated clean K/V shard serves every local block query whose prefix
includes that shard. In backward, corrupted K/V gradients remain on block
owners and clean K/V gradients return to their shard owners. Cross-rank
attention communication therefore contains only shared clean K/V and their
gradients. CSBP removes the length-$L$ corrupted component from this
communication in both forward and backward.

\paragraph{Memory.}
Each CSBP rank stores one length-$L/P$ clean shard and the activations for its
assigned corrupted blocks. With balanced block assignments, each rank stores
approximately $L/P$ clean tokens and $L/P$ corrupted tokens, giving
approximately $2L/P$ token activations and matching conventional CP to first
order. Pure BP can store a clean prefix approaching length $L$ on ranks
assigned later blocks. CSBP reduces this clean-prefix component to
$L/P$, providing up to a $P\times$ reduction for a rank whose longest required
prefix approaches $L$. CSBP therefore retains CP's activation scaling
without pure BP's overlapping-prefix replication.
